\Needspace{14\baselineskip}
\section{Conclusion}
\label{sec:conclusion}

This systematic review mapped 28 studies on knowledge persistence in
AI-based coding agents, all published from 2024 onwards.
The field is recent and predominantly composed of preprints (82\%) but
already shows convergence around recurring architectural and evaluation
patterns.
Memory persistence produces gains in every tested scenario, but with
diminishing returns, documented degradation risks, and persistent gaps
in cost reporting.

The five secondary questions received answers with differing degrees of
certainty.
For the taxonomy (SQ1), experience banks predominate (46\%), the LLM
acts as the controller in 54\% of the systems, and hybrid temporal
scope is the most frequent (58\%).
On evaluation (SQ2), SWE-bench Verified concentrates 67\% of studies
with a benchmark, though comparability is limited by differences in
model, framework, and experimental conditions.
Performance gains (SQ3) are universally positive across the 19
controlled comparisons, with a median of 6.8~pp and an amplitude of
+2.2 to +39.0~pp, and diminishing returns for baselines above 65\%.
For cost (SQ4), 43\% of the studies report no cost data, and the
memory impact is bimodal: systems that replace context reduce tokens,
while systems that add context raise cost by approximately 5 to 21\%.
On complexity and effectiveness (SQ5), 32\% of the studies report
memory-induced degradation, with retrieval noise as the main mechanism,
and no evidence exists that more complex architectures yield
systematically larger gains.

None of the 28 studies compares persistence architectures directly.
This absence, combined with cost underreporting and the uncertainty
about the role of architectural complexity, prevents informed decisions
about which approach to adopt in each context.
The most critical gap --- the absence of direct comparisons --- calls
for a comparative experiment with at least three representative
architectures, under the same benchmark, the same model, and with
integrated per-task cost analysis.
